%========================================
% MatinBook — Stage 02: Font System Test
% Version: v1.1
% Purpose:
%   Validate the entire font system of MatinBook v1.1:
%     1. Persian body font (XB Niloofar)
%     2. Latin body font (Times New Roman)
%     3. Math font (XITS Math)
%     4. Monospace font (DejaVu Sans Mono)
%     5. Code font (JetBrains Mono) with ligatures
%     6. Mixed Persian/Latin paragraphs
%     7. Multilingual quotations
%     8. Font comparison table
%
% Compilation:
%   xelatex -shell-escape stage02-fonts.tex
%   xelatex -shell-escape stage02-fonts.tex
%   (2 passes required for cover + cross-references)
%
% Expected outcome:
%   A professionally typeset PDF demonstrating that every font
%   family used by MatinBook renders correctly, with proper
%   kerning, ligatures, and mixed-direction handling.
%========================================

%------------------------------------------------------------------------
% Search Path (for tests directory)
%------------------------------------------------------------------------
\makeatletter
\def\input@path{{../}}
\makeatother

%------------------------------------------------------------------------
% Document Class
%------------------------------------------------------------------------
\documentclass{matinbook}

%------------------------------------------------------------------------
% Book Metadata
%------------------------------------------------------------------------
\title{آزمون سیستم فونت‌ها}
\author{تیم توسعه MatinBook}
\date{مهر ۱۴۰۵}

\booktitle{آزمون سیستم فونت‌ها}

\renewcommand{\repository}{github.com/matinmahmoudi-cs/matinbook}

%========================================================================
% DOCUMENT
%========================================================================
\begin{document}

%------------------------------------------------------------------------
% FRONT COVER
%------------------------------------------------------------------------
\makecover
    {آزمون سیستم فونت‌ها}
    {ارزیابی خانواده‌های فونت در MatinBook}
    {تیم توسعه MatinBook}
    {مهر ۱۴۰۵}

%------------------------------------------------------------------------
% FRONT MATTER (symmetric margins)
%------------------------------------------------------------------------
\frontmatter
\frontmattergeometry

% Title page
\maketitle

% Copyright / Colophon
\thispagestyle{empty}
\vfill
\begin{center}
    {\large\textbf{شناسنامه آزمون}}

    \vspace{0.8cm}

    \begin{tabular}{rl}
        عنوان: & آزمون سیستم فونت‌ها \\
        نسخه: & \matinbookversion \\
        تاریخ: & \matinbookdate \\
        موتور: & \lr{XeLaTeX} \\
        مجوز: & \lr{MIT} \\
    \end{tabular}

    \vspace{0.8cm}

    {\small این سند بخشی از مجموعه آزمون‌های یکپارچگی \lr{MatinBook} است.}
\end{center}
\clearpage

% Preface
\chapter*{پیشگفتار}
\addcontentsline{toc}{chapter}{پیشگفتار}

این سند، دومین آزمون از مجموعه‌ی آزمون‌های یکپارچگی \lr{MatinBook}
نسخه‌ی \lr{\matinbookversion} است. هدف آن، اعتبارسنجی کامل سیستم
فونت‌های کلاس در چهار حوزه‌ی متمایز است: متن فارسی، متن لاتین،
ریاضیات و کد.

در یک کتاب فنی فارسی، هماهنگی میان این چهار حوزه، تعیین‌کننده‌ی
کیفیت نهایی حروف‌چینی است. فونت فارسی باید در کنار فونت لاتین
تعادل بصری داشته باشد؛ فونت ریاضی باید نمادهای پیچیده را بدون
افت کیفیت نمایش دهد؛ و فونت کد باید لیگچرهای برنامه‌نویسی را
به‌درستی پشتیبانی کند.

هر بخش از این آزمون، با مثال‌های عملی و ارجاعات متقابل همراه است
تا خواننده بتواند درستی رفتار هر فونت را در بافت واقعی یک کتاب
فنی ارزیابی کند.

\clearpage

% Table of Contents
\tableofcontents
\clearpage

% List of Figures
\listoffigures
\clearpage

% List of Tables
\listoftables
\clearpage

%------------------------------------------------------------------------
% MAIN MATTER (wide margin for notes)
%------------------------------------------------------------------------
\mainmatter
\mainmattergeometry

%========================================================================
% CHAPTER 1 — PERSIAN FONT
%========================================================================
\chapter{فونت فارسی}
\label{chap:font-persian}

\section{فلسفه‌ی انتخاب فونت}
\label{sec:font-persian-philosophy}

فونت پیش‌فرض متن فارسی در \lr{MatinBook}، \lr{XB Niloofar} است
که با مقیاس \pnum{۱.۲} بارگذاری می‌شود. این فونت به‌دلیل پوشش کامل
حروف فارسی، خوانایی بالا در متون فنی، و تعادل بصری با فونت لاتین
\lr{Times New Roman} انتخاب شده است.

مقیاس \pnum{۱.۲} به این معناست که اندازه‌ی فونت فارسی، ۲۰٪ بزرگ‌تر از
اندازه‌ی پایه‌ی سند (\pnum{۱۰pt}) است. این تنظیم، نتیجه‌ی سال‌ها تجربه‌ی
نشر فارسی است: فونت‌های فارسی به‌دلیل ساختار منحنی‌دارشان، در
اندازه‌های کوچک، خوانایی کمتری نسبت به فونت‌های لاتین دارند.

\mnote{مقیاس \pnum{۱.۲} برای فونت فارسی، استاندارد نشر فنی ایران است.}

\section{خوانایی در متن طولانی}
\label{sec:font-persian-readability}

این یک پاراگراف طولانی فارسی برای ارزیابی خوانایی فونت
\lr{XB Niloofar} است. فونت فارسی باید به گونه‌ای طراحی شده باشد
که خواننده بتواند ساعت‌های طولانی بدون خستگی چشم مطالعه کند.
فاصله بین حروف، اندازه نقاط، و شکل منحنی‌ها همگی در خوانایی
متن تأثیرگذار هستند.

طراحی حروف در \lr{XB Niloofar} بر پایه‌ی اصول سنتی خط نسخ است،
اما با تنظیمات مدرن برای نمایش دیجیتال. این ترکیب، هم حس
آشنایی برای خواننده‌ی فارسی‌زبان ایجاد می‌کند و هم در چاپ
دیجیتال، کیفیت بالایی دارد.

\section{حروف خاص فارسی}
\label{sec:font-persian-special}

حروف فارسی شامل ۳۲ حرف اصلی است. چهار حرف در فارسی وجود دارند
که در عربی معادل ندارند و به همین دلیل، پشتیبانی فونت از آن‌ها
معیار مهمی برای کیفیت است:

\begin{itemize}
    \item \textbf{«گ»} — مانند: گفتگو، بزرگ، برگ، گرگ
    \item \textbf{«چ»} — مانند: چرا، هیچ، کاغذ، چراغ
    \item \textbf{«پ»} — مانند: پدر، پاسخ، پرونده، پژواک
    \item \textbf{«ژ»} — مانند: ژرف، مژده، پژوهش، ژاله
\end{itemize}

\section{اعداد فارسی}
\label{sec:font-persian-numbers}

اعداد فارسی باید در تمام حالت‌ها به‌درستی نمایش داده شوند.
این شامل اعداد ساده، دو رقمی، چند رقمی، اعشاری و درصدی است:

\begin{itemize}
    \item اعداد ساده: ۰ ۱ ۲ ۳ ۴ ۵ ۶ ۷ ۸ ۹
    \item اعداد دو رقمی: ۱۰ ۲۵ ۴۸ ۷۳ ۹۹
    \item اعداد چند رقمی: ۱۴۰۴ ۲۰۲۶ ۱۲۳۴۵۶ ۹۸۷۶۵۴۳۲۱۰
    \item اعداد اعشاری: \pnum{۳/۱۴}، \pnum{۲/۷۱۸}، \pnum{۱/۶۱۸}
    \item درصد: ۲۵٪ ۵۰٪ ۷۵٪ ۱۰۰٪
\end{itemize}

\section{نویسه‌های نگارشی}
\label{sec:font-persian-punctuation}

نویسه‌های ویژه‌ی فارسی و نشانه‌گذاری:

\begin{itemize}
    \item گیومه فارسی: «نقل قول مستقیم»
    \item گیومه تکی: ‹نقل قول درونی›
    \item پرانتز: (متن داخل پرانتز)
    \item کروشه: [متن داخل کروشه]
    \item آکولاد: \{متن داخل آکولاد\}
    \item ممیز فارسی: \pnum{۳/۱۴۱۵۹۲}
    \item جداکننده هزارگان: \pnum{۱٬۰۰۰٬۰۰۰}
\end{itemize}

\section{نمونه‌ی متن ادبی}
\label{sec:font-persian-poetry}

برای ارزیابی زیبایی‌شناختی فونت در متن‌های ادبی، دو بیت از
شعر کلاسیک فارسی در ادامه آمده است:

\begin{flushright}
    \textbf{توانا بود هر که دانا بود}\\
    \textit{ز دانش دل پیر برنا بود}

    \vspace{0.5cm}
    به نام خداوند جان و خرد\\
    کزین برتر اندیشه برنگذرد
\end{flushright}

%========================================================================
% CHAPTER 2 — LATIN FONT
%========================================================================
\chapter{فونت لاتین}
\label{chap:font-latin}

\section{فونت Times New Roman}
\label{sec:font-latin-times}

\begin{latin}
This chapter demonstrates the Latin font \textbf{Times New Roman},
which is the default Latin font in MatinBook. This serif font is
widely used in academic and technical publications due to its
excellent readability and professional appearance.
\end{latin}

\section{آزمون پانگرام}
\label{sec:font-latin-pangram}

\begin{latin}
A pangram is a sentence that contains every letter of the alphabet
at least once. The following sentences test the font's handling of
all 26 letters:

The quick brown fox jumps over the lazy dog.

Pack my box with five dozen liquor jugs.

How vexingly quick daft zebras jump!
\end{latin}

\section{لیگچرها}
\label{sec:font-latin-ligatures}

\begin{latin}
Common ligatures in English text:

\begin{itemize}
    \item fi --- finance, define, profile
    \item fl --- flower, reflect, inflate
    \item ff --- offer, different, difficult
    \item ffi --- efficient, coefficient
    \item ffl --- shuffle, waffle
\end{itemize}
\end{latin}

\section{نویسه‌های ویژه}
\label{sec:font-latin-special}

\begin{latin}
\begin{itemize}
    \item Quotes: ``double'' and `single'
    \item Dashes: -- (en-dash) and --- (em-dash)
    \item Symbols: @ \# \$ \% \^{} \& * \_ + = | \~{} < > /
    \item Brackets: [square], \{curly\}, (round)
\end{itemize}
\end{latin}

\section{متن علمی لاتین}
\label{sec:font-latin-scientific}

\begin{latin}
\textbf{Computer Science Terms.}
The time complexity of the algorithm is \(O(n \log n)\) in the
average case and \(O(n^2)\) in the worst case. We use a
\texttt{HashMap} data structure for \(O(1)\) lookup time.

\textbf{Mathematics Terms.}
The Gaussian integral is defined as:
\[\int_{-\infty}^{\infty} e^{-x^2}\,dx = \sqrt{\pi}\]
The Riemann zeta function for \(s > 1\) is:
\[\zeta(s) = \sum_{n=1}^{\infty} \frac{1}{n^s}\]
\end{latin}

%========================================================================
% CHAPTER 3 — MATH FONT
%========================================================================
\chapter{فونت ریاضی}
\label{chap:font-math}

\section{الفبای یونانی}
\label{sec:font-math-greek}

جدول \cref{tab:greek-alphabet} الفبای یونانی را با نام
\lr{LaTeX} هر نماد نشان می‌دهد. این نمادها در فونت \lr{XITS Math}
با کیفیت بالا رندر می‌شوند.

\begin{table}[h]
\centering
\caption{الفبای یونانی در فونت \lr{XITS Math}}
\label{tab:greek-alphabet}
\begin{tabular}{cc@{\hspace{1.5cm}}cc}
\toprule
\textbf{Uppercase} & \textbf{Command} & \textbf{Lowercase} & \textbf{Command} \\
\midrule
\(\Gamma\) & \texttt{\textbackslash Gamma} & \(\alpha\) & \texttt{\textbackslash alpha} \\
\(\Delta\) & \texttt{\textbackslash Delta} & \(\beta\)  & \texttt{\textbackslash beta} \\
\(\Theta\) & \texttt{\textbackslash Theta} & \(\gamma\) & \texttt{\textbackslash gamma} \\
\(\Lambda\) & \texttt{\textbackslash Lambda} & \(\delta\) & \texttt{\textbackslash delta} \\
\(\Pi\)    & \texttt{\textbackslash Pi}    & \(\epsilon\) & \texttt{\textbackslash epsilon} \\
\(\Sigma\) & \texttt{\textbackslash Sigma} & \(\zeta\)  & \texttt{\textbackslash zeta} \\
\(\Phi\)   & \texttt{\textbackslash Phi}   & \(\theta\) & \texttt{\textbackslash theta} \\
\(\Omega\) & \texttt{\textbackslash Omega} & \(\lambda\) & \texttt{\textbackslash lambda} \\
\bottomrule
\end{tabular}
\end{table}

\section{عملگرهای ریاضی}
\label{sec:font-math-operators}

نمایش نمادهای اصلی ریاضی در فونت \lr{XITS Math}:

\begin{itemize}
    \item جمع و ضرب: $\sum_{i=1}^{n} i = \frac{n(n+1)}{2}$،
          $\prod_{k=1}^{n} k = n!$
    \item انتگرال: $\int_{0}^{\infty} e^{-x} \,dx = 1$،
          $\iint_{D} f(x,y) \,dA$
    \item حد: $\lim_{x \to \infty} (1 + \frac{1}{x})^{x} = e$
    \item مشتق: $\frac{d}{dx}(x^{n}) = nx^{n-1}$،
          $\frac{\partial^{2}f}{\partial x^{2}}$
    \item رادیکال: $\sqrt{x^{2}+y^{2}}$،
          $\sqrt[3]{8} = 2$
\end{itemize}

\section{ماتریس و بردار}
\label{sec:font-math-matrices}

نمایش ماتریس و بردار با فونت ریاضی:

\[
A_{3 \times 3} =
\begin{bmatrix}
    a_{11} & a_{12} & a_{13} \\
    a_{21} & a_{22} & a_{23} \\
    a_{31} & a_{32} & a_{33}
\end{bmatrix}, \quad
\vec{v} = \begin{bmatrix} v_{1} \\ v_{2} \\ v_{3} \end{bmatrix}
\]

دترمینان ماتریس $A$:

\[
\det(A) = a_{11}(a_{22}a_{33} - a_{23}a_{32})
        - a_{12}(a_{21}a_{33} - a_{23}a_{31})
        + a_{13}(a_{21}a_{32} - a_{22}a_{31})
\]

\section{معادلات معروف}
\label{sec:font-math-famous}

چهار معادله‌ی مشهور ریاضی که برای ارزیابی کیفیت فونت
ریاضی در سطوح مختلف پیچیدگی انتخاب شده‌اند:

\begin{align}
    e^{i\pi} + 1 &= 0 \label{eq:euler} \\
    \int_{-\infty}^{\infty} e^{-x^{2}} \,dx &= \sqrt{\pi} \label{eq:gaussian} \\
    \sum_{n=1}^{\infty} \frac{1}{n^{2}} &= \frac{\pi^{2}}{6} \label{eq:basel} \\
    \zeta(s) &= \prod_{p \text{ prime}} \frac{1}{1 - p^{-s}} \label{eq:euler-product}
\end{align}

معادله‌ی \cref{eq:euler}، که به‌عنوان «زیباترین معادله‌ی ریاضی»
شناخته می‌شود، چهار ثابت بنیادین ریاضی را در یک رابطه‌ی ساده
به هم پیوند می‌دهد. معادله‌ی \cref{eq:gaussian}، انتگرال گاوسی
است که در نظریه‌ی احتمال نقش محوری دارد. معادله‌ی
\cref{eq:basel}، مسئله‌ی بازل است که اویلر در قرن هجدهم حل کرد.
و معادله‌ی \cref{eq:euler-product}، رابطه‌ی میان تابع زتای
ریمان و اعداد اول را نشان می‌دهد.

%========================================================================
% CHAPTER 4 — CODE FONT
%========================================================================
\chapter{فونت کد}
\label{chap:font-code}

\section{فونت JetBrains Mono}
\label{sec:font-code-jetbrains}

\lr{JetBrains Mono} یک فونت \lr{monospace} مدرن است که مخصوص
برنامه‌نویسی طراحی شده. این فونت دارای لیگچرهای مخصوص کدنویسی
است که خوانایی کد را افزایش می‌دهند. لیگچرها ترکیب‌های ویژه‌ای
از دو یا چند نویسه هستند که به‌صورت یک نماد واحد نمایش داده
می‌شوند.

\begin{latin}
\begin{minted}{text}
    Comparison:  == != <= >= -> => :: <| |>
    Arrows:      <- -> --> <-- -->>
    Logic:       && || !!
    Math:        ++ -- ** // /= =/=
\end{minted}
\end{latin}

\section{نمونه کد پایتون}
\label{sec:font-code-python}

نمونه‌ی \cref{code:greeting} یک تابع پایتون با \lr{type hints}
و \lr{docstring} را نشان می‌دهد. این نمونه برای ارزیابی
خوانایی فونت کد در متن‌های واقعی طراحی شده است.

\begin{latin}
\begin{minted}{python}
def greet(name: str, age: int = 0) -> str:
    """Return a greeting message."""
    if age > 0:
        return f"Hello {name}, you are {age} years old!"
    return f"Hello {name}!"

# Test the function
names = ["Ali", "Sara", "Reza"]
for i, name in enumerate(names, start=1):
    print(f"{i}. {greet(name, 20 + i)}")
\end{minted}
\end{latin}
\codecaption{تابع greeting با type hints}
\label{code:greeting}

\section{لیگچرهای کد}
\label{sec:font-code-ligatures}

برای ارزیابی پشتیبانی از لیگچرها، کد زیر ترکیب‌های مختلف را
نمایش می‌دهد:

\begin{latin}
\begin{minted}{cpp}
template <typename T>
constexpr auto compute(T x, T y) -> T {
    auto result = (x != 0) ? x * y : y / x;
    return result <= 100 ? result : 100;
}
\end{minted}
\end{latin}
\codecaption{نمونه کد C++ با لیگچرها}
\label{code:cpp-ligatures}

نمونه‌ی \cref{code:cpp-ligatures} شامل لیگچرهای \lr{<=}،
\lr{!=}، \lr{->}، \lr{?} و \lr{:} است که همگی باید در فونت
\lr{JetBrains Mono} به‌درستی رندر شوند.

%========================================================================
% CHAPTER 5 — MULTILINGUAL TEXT
%========================================================================
\chapter{متن چندزبانه}
\label{chap:multilingual}

\section{ترکیب فارسی و لاتین}
\label{sec:multilingual-mixing}

در متن‌های فنی فارسی، ترکیب متن فارسی و لاتین امری اجتناب‌ناپذیر
است. \lr{MatinBook} با استفاده از بسته‌ی \lr{xepersian}، این
ترکیب را به‌صورت خودکار مدیریت می‌کند: جهت متن بر اساس زبان
جاری تعیین می‌شود و فونت مناسب برای هر بخش انتخاب می‌شود.

این یک پاراگراف ترکیبی است که در آن از واژگان انگلیسی
مانند \lr{function}، \lr{variable}، \lr{array}، \lr{loop}،
\lr{conditional} و \lr{recursion} استفاده شده است.
همچنین ممکن است اعدادی مانند \pnum{۳/۱۴۱۵۹} یا \pnum{۲/۷۱۸۲۸}
در متن فارسی ظاهر شوند. عبارات ریاضی مانند $f(x) = x^{2}$
نیز در متن فارسی قرار می‌گیرند.

نام دانشمندان خارجی: \lr{Isaac Newton}، \lr{Albert Einstein}،
\lr{Alan Turing}، \lr{Donald Knuth}، \lr{Ada Lovelace}

\section{نقل قول چندزبانه}
\label{sec:multilingual-quotes}

نقل قول‌های انگلیسی در متن فارسی، نیازمند مدیریت دقیق جهت
و فونت هستند. نقل قول زیر، یکی از معروف‌ترین جملات در
مهندسی نرم‌افزار است:

\begin{latin}
\begin{quote}
    ``Any fool can write code that a computer can understand.
    Good programmers write code that humans can understand.''

    \hfill --- \textit{Martin Fowler}
\end{quote}
\end{latin}

ترجمه‌ی این نقل قول در ادامه آمده است:

\begin{flushright}
    «هر احمقی می‌تواند کدی بنویسد که کامپیوتر بفهمد.
    برنامه‌نویسان خوب کدی می‌نویسند که انسان‌ها بفهمند.»

    \hfill --- مارتین فاولر
\end{flushright}

%========================================================================
% CHAPTER 6 — SUMMARY
%========================================================================
\chapter{خلاصه‌ی نتایج}
\label{chap:summary}

\section{وضعیت فونت‌ها}
\label{sec:summary-fonts}

جدول \cref{tab:fonts} وضعیت چهار خانواده‌ی فونت \lr{MatinBook}
را نشان می‌دهد. هر خانواده برای نقش مشخصی در کتاب طراحی شده
و هماهنگی میان آن‌ها، کیفیت نهایی حروف‌چینی را تعیین می‌کند.

\begin{table}[h]
\centering
\caption{وضعیت خانواده‌های فونت \lr{MatinBook}}
\label{tab:fonts}
\begin{tabular}{@{}l l C{3cm} C{4cm}@{}}
\toprule
\textbf{نوع} & \textbf{فونت} & \textbf{وضعیت} & \textbf{توضیحات} \\
\midrule
فارسی & \lr{XB Niloofar} & \textcolor{matingreen}{فعال} & مقیاس \pnum{۱.۲}، خوانایی بالا \\
لاتین & \lr{Times New Roman} & \textcolor{matingreen}{فعال} & سریف، استاندارد آکادمیک \\
ریاضی & \lr{XITS Math} & \textcolor{matingreen}{فعال} & سازگار با \lr{unicode-math} \\
کد & \lr{JetBrains Mono} & \textcolor{matingreen}{فعال} & لیگچرهای برنامه‌نویسی \\
\bottomrule
\end{tabular}
\end{table}

\section{معیارهای ارزیابی}
\label{sec:summary-criteria}

در این آزمون، هر خانواده‌ی فونت بر اساس چهار معیار ارزیابی شد:

\begin{enumerate}
    \item \textbf{خوانایی}: در متن‌های طولانی و در اندازه‌های کوچک.
    \item \textbf{پوشش نویسه}: شامل حروف خاص، اعداد، و نشانه‌گذاری.
    \item \textbf{هماهنگی}: تعادل بصری با سایر فونت‌های سند.
    \item \textbf{پشتیبانی فنی}: لیگچرها، اعداد کسری، و نمادهای ویژه.
\end{enumerate}

\section{نتیجه‌گیری}
\label{sec:summary-conclusion}

\textbf{تمام چهار خانواده‌ی فونت \lr{MatinBook} نسخه‌ی
\lr{\matinbookversion} با موفقیت بارگذاری و آزمایش شدند.
نمایش فارسی، لاتین، ریاضی و کد همگی بدون نقص هستند و
سیستم فونت کلاس، آماده‌ی استفاده در پروژه‌های واقعی است.}

%------------------------------------------------------------------------
% BACK MATTER (symmetric margins)
%------------------------------------------------------------------------
\backmatter
\frontmattergeometry

% Bibliography (empty placeholder)
\printbibliography[title={منابع و مراجع}]

% Index
\printindex

%------------------------------------------------------------------------
% BACK COVER
%------------------------------------------------------------------------
\authorbio{%
    تیم توسعه‌ی \lr{MatinBook}، گروهی از پژوهشگران و برنامه‌نویسان
    فارسی‌زبان است که با هدف ارائه‌ی یک راه‌حل حرفه‌ای برای
    حروف‌چینی کتاب‌های فنی فارسی فعالیت می‌کند. این پروژه حاصل
    سال‌ها تجربه در حوزه‌ی نشر دیجیتال و برنامه‌نویسی است.
}

\makebackcover{%
    این سند، دومین آزمون از مجموعه‌ی آزمون‌های یکپارچگی
    \lr{MatinBook} نسخه‌ی \lr{\matinbookversion} است.

    \vspace{0.5cm}

    \textbf{زیرسیستم‌های آزموده‌شده:}
    \begin{itemize}
        \item فونت فارسی \lr{XB Niloofar}
        \item فونت لاتین \lr{Times New Roman}
        \item فونت ریاضی \lr{XITS Math}
        \item فونت مونواسپیس \lr{DejaVu Sans Mono}
        \item فونت کد \lr{JetBrains Mono}
        \item متن چندزبانه
    \end{itemize}
}

\end{document}